\section{No usable method finds one}\label{sec:discover}

\Cref{thm:k1-certify} rules out certifying a final theory. A weaker goal remains: a
method might converge on a final theory without ever announcing that it has. This
section shows that no usable method is guaranteed, even with any fixed probability, to
settle on a final theory, and that whatever success a method has is bought by a bias
the record does not supply.

Fix the length-lexicographic enumeration $x_0,x_1,x_2,\dots$ of $\Str$, a computable
bijection $\N\to\Str$, and identify a total $f:\Str\to\Str$ with the sequence
$f(x_0),f(x_1),\dots$. For a finite sequence $\sigma=(y_0,\dots,y_{n-1})$ write
$\rec_\sigma=\big((x_0,y_0),\dots,(x_{n-1},y_{n-1})\big)$, write $\sigma\tau$ for
concatenation, $\sigma\sqsubseteq\tau$ for the prefix order, and
$f\restriction n=(f(x_0),\dots,f(x_{n-1}))$. For a record $\rec$ already in hand,
$\rec\frown\rec_\tau$ is the record $\rec$ followed by the observations of $\rec_\tau$;
the presentation of $f$ after $\rec$ is the chain
$\rec\frown\rec_{f\restriction 0}\sqsubseteq\rec\frown\rec_{f\restriction 1}
\sqsubseteq\cdots$ (\cref{def:learner}). Fix two distinct symbols $0,1\in\Sig$.

Any finite functional record $\rec$ determines a finite sequence $\sigma_\rec$: take
the least $n$ with $\dom(\rec)\subseteq\{x_0,\dots,x_{n-1}\}$ and set the $i$th entry to
the outcome $\rec$ records at $x_i$, or to $0$ if it records none. Every total $f$
extending $\sigma_\rec$ is consistent with $\rec$.

\begin{theorem}[T2: Non-identifiability from any record; after \citet{putnam1963,gold1967}]\label{thm:ud-gold}
Let $L$ be a computable method (\cref{def:learner}), possibly partial, and $\rec$ a
finite functional record. There is a total computable world $f$,
consistent with $\rec$, that $L$ does not identify in the limit after $\rec$.
\end{theorem}

\begin{proof}
We build finite sequences $\sigma_0\sqsubseteq\sigma_1\sqsubseteq\cdots$ with
$\sigma_0=\sigma_\rec$. Given $\sigma_s$, search, dovetailing over $k$ and computation
time, for a $k$ such that $h=L(\rec\frown\rec_{\sigma_s0^k})$ is defined and
\[
  \ph{h}\bigl(x_{|\sigma_s|+k}\bigr)\conv = 0,
\]
that is, the current conjecture predicts that the record continues with $0$. If the
search succeeds, set $\sigma_{s+1}=\sigma_s0^k1$.

Suppose the search at some stage $s$ never succeeds. Then $L$ does not identify
$g=\sigma_s0^\infty$ after $\rec$. If it did, it would be defined on every record of
the presentation of $g$ after $\rec$, and from some point on its conjecture would be a
fixed index for $g$; that index outputs $0$ at every argument beyond $\sigma_s$, so the
search would succeed. The world $g$ is total computable and extends $\sigma_\rec$, so it
is consistent with $\rec$.

Otherwise every stage succeeds, and $f=\bigcup_s\sigma_s$ is total computable and
extends $\sigma_\rec$, so it is consistent with $\rec$. The conjecture $h_s$ found at
stage $s$ is the output of $L$ on a record of the presentation of $f$ after $\rec$,
since $\sigma_s0^k\sqsubseteq f$, and it predicts $0$ where $f$ has $1$, so
$\ph{h_s}\ne f$. These records grow with $s$. If $L$ converged on $f$ to some $h^*$, then
$h_s=h^*$ for all large $s$, so $\ph{h^*}\ne f$.
\end{proof}

A method that does well on the simplest worlds (here, those that are eventually
constant) for that very reason never settles on the right theory in another computable
world, and that world can be chosen to agree with any record already in hand. The world
is not exotic. The construction is uniform in $L$ and $\rec$. If every stage succeeds,
$\K(f)\le\K(L,\rec)+O(1)\le\K(L)+\K(\rec)+O(1)$, where $\K(L)$ is the complexity of an
index for $L$. If stage $s$ stalls, the world is eventually constant and
$\K(g)\le\K(L,\rec)+\K(s)+O(1)$. Either way its complexity is set by the method, the
record and at most one stage number, and its records are compressible, not atypical or
non-stochastic in the sense of algorithmic statistics \citep{gacs2001}.

\begin{remark}[Randomized methods]\label{ud:rem-randomized}
A method may also flip a coin of computable bias. The theorem covers such methods in
the following form. A probabilistic method
that identifies every computable world with probability greater than $1/2$ can be
replaced by a deterministic computable method that identifies the same worlds
\citep[Thm.~3.21 with $n=1$]{pitt1989}, so by \cref{thm:ud-gold} no probabilistic method identifies every
computable world with probability greater than $1/2$. Success with a smaller probability, bounded below by
a fixed $p>0$ for every world, corresponds to identification by a finite team of
deterministic methods, and the class of all computable worlds is not identifiable by
any finite team either \citetext{\citealp[Thm.~3.36]{pitt1989}; \citealp[Thm.~3.1]{smith1982}; \citealp{pitt-smith1988}}. Both facts hold
relative to any record: the class of computable worlds extending $\sigma_\rec$ is not
finite-team identifiable, since a team for it would yield one for all computable worlds
via $g\mapsto\sigma_\rec g$ and the corresponding computable shift of indices. So, given
any record, no usable randomized method that receives the presentation identifies every
computable world consistent with it with any fixed positive probability. This is a
uniform statement: it does not say that the method fails almost surely in some single
world, and a randomized method can identify every computable world with \emph{some}
positive probability, for instance by guessing an index at random and keeping it.
\end{remark}

\begin{remark}[Active methods]\label{rem:active}
A method may also choose which settings to observe. An \emph{active} method reads the
record so far, and possibly random bits, and outputs a conjecture and a next setting to
query; the world answers the query, and the answer extends the record. Active methods
gain nothing here. Given a computable active method $M$ and a record $\rec$, let $L_M$
be the method that receives the presentation after $\rec$ and, on the record
$\rec\frown\rec_{f\restriction n}$, simulates $M$ from $\rec$ for $n$ computation steps
with the same random bits, answering each query from the record and halting the
simulation at the first query the record does not yet answer; $L_M$ outputs the last
conjecture of $M$ reached. $L_M$ is computable, and randomized exactly when $M$ is.
Every setting is eventually answered, so every step of a run of $M$ in which each
step halts is eventually reached; if on given random bits $M$ identifies $f$, its
conjectures are eventually constant at an index of $f$, and so are those of $L_M$.
Hence $L_M$ identifies every world that $M$ identifies, with at least the same
probability. Applied to $L_M$, \cref{thm:ud-gold} and \cref{ud:rem-randomized} give,
for every record, a computable world consistent with it that $M$ does not identify, or,
for randomized $M$ and every $p>0$, one that $M$ identifies with probability less than
$p$.
\end{remark}

\begin{corollary}[No record rules out failure]\label{ud:cor-every-method-fails}
Let $L$ be a usable method, possibly partial, and $\rec$ a finite functional record that
$L$ has observed; $L$ may receive the presentation or choose its own queries. If $L$
is deterministic, some computable world consistent with $\rec$ \emph{defeats} it: in
that world, $L$ never settles on the right theory. If $L$ is randomized, then for every
$p>0$ there is a computable world consistent with $\rec$ that $L$ identifies with
probability less than $p$. In either case no
record tells $L$'s user that the world is not such a world.
\end{corollary}

\begin{proof}
By \cref{prem:effective} $L$ is computable or a computable probabilistic method;
apply \cref{thm:ud-gold} or \cref{ud:rem-randomized}, and for active methods
\cref{rem:active}, to $\rec$.
\end{proof}

\paragraph{The limit does not depend on the level of computing power.}
\Cref{thm:ud-gold} is not an artifact of ordinary computers. Fix any set $A\subseteq\N$
as an oracle, and let $(\ph{e}^A)_{e\in\N}$ be an acceptable numbering of the functions
partial computable relative to $A$, which exists by the relativized universal-function
and $s$-$m$-$n$ theorems \citep[ch.~9]{rogers1967}; the proof below uses only that the universal
function is $A$-computable. Call a model at level $A$ an index $e$ of
$\ph{e}^A$, and a world at level $A$ a total $A$-computable function; identification
at level $A$ is defined as in \cref{def:learner}, with $\ph{e}^A$ in place of
$\ph{e}$.

\begin{theorem}[T2 at every level]\label{thm:ud-gold-rel}
Let $A\subseteq\N$, let $L$ be an $A$-computable method, possibly partial, that outputs
models at level $A$, that is, indices in $(\ph{e}^A)$ rather than ordinary indices, and
let $\rec$ be a finite functional record. There is a total $A$-computable world $f$,
consistent with $\rec$, that $L$ does not identify at level $A$ after $\rec$.
\end{theorem}

\begin{proof}
The proof of \cref{thm:ud-gold} goes through with $\ph{}^A$ in place of $\ph{}$. Given
$\sigma_s$ (with $\sigma_0=\sigma_\rec$), search for $k$ such that
$h=L(\rec\frown\rec_{\sigma_s0^k})$ is defined and $\ph{h}^A(x_{|\sigma_s|+k})\conv=0$.
The computations of $L$ and of $\ph{h}^A$ can be simulated with access to $A$, so the
search, dovetailed over $k$ and computation time, is $A$-effective. If it never
succeeds at stage $s$, then $L$ does not identify $\sigma_s0^\infty$ at level $A$, for
the same reason as before, and that world is computable, hence $A$-computable.
Otherwise set $\sigma_{s+1}=\sigma_s0^k1$. Then $f=\bigcup_s\sigma_s$ is total, extends
$\sigma_\rec$, and is $A$-computable, since every stage is effective relative to $A$.
Each conjecture $h_s$ found at stage $s$ predicts $0$ where $f$ has $1$, so if $L$
converged on $f$ to some $h^*$ then $\ph{h^*}^A\ne f$.
\end{proof}

The relativized theorem needs the worlds to be at the method's level: a method strictly
more powerful than the worlds it faces can succeed, as the ideal-agent concession below
shows. For deterministic methods, C2 therefore needs only the following premise.

\begingroup
\renewcommand{\thepremise}{P2$^{*}$}%
\renewcommand{\theHpremise}{P2star}%
\addtocounter{premise}{-1}%
\begin{premise}[Methods and worlds share a level]\label{rem:p2star}
For every usable method $L$ there is a set $A\subseteq\N$ such that $L$ is $A$-computable, its outputs are read as models at level $A$, and its guarantee must cover every total $A$-computable world consistent with the record: every such world is a possible world. A model at level $A$ is final for a world $f$ if it computes $f$ relative to $A$.
\end{premise}
\endgroup

\noindent\emph{Refuted by:} a usable method more powerful than every world its guarantee
must cover, for instance a usable oracle whose power the theories it chooses among
cannot share.

\Cref{prem:effective}, together with the guarantee of \cref{def:learner}, which covers
every total computable world consistent with the record, is the case $A=\emptyset$ of
P2$^{*}$. Under P2$^{*}$, \cref{thm:ud-gold-rel} applies at the method's own level, so
a usable hypercomputer of oracle strength, whose power the theories share, does not
escape C2. Hypercomputers not equivalent to a single oracle, such as infinite-time
Turing machines, would need their own argument. The
reductions of \citet{pitt1989} behind \cref{ud:rem-randomized} are uniform simulations
and should relativize, but we have not verified this, so P2$^{*}$ serves for
deterministic methods only.

\begin{lemma}[L: Every consistent method is a bias]\label{lem:bias}
Let $L$ be a method that outputs, on each functional record $\rec$, a model in
$\Cons(\rec)$.
\begin{enumerate}[label=(\roman*)]
  \item $L$ is a record-indexed constraint (\cref{def:constraint}): with
        $c_\rec(e)=0$ if $e=L(\rec)$ and $c_\rec(e)=1$ otherwise,
        $\Sel_c(\rec)=\{L(\rec)\}$. For every functional record $\rec$ there is a
        method $L'$ that outputs consistent models, is computable if $L$ is, and
        selects differently from $L$ at $\rec$.
  \item For a finite functional record $\rec$, let $K_L$ be the class of total
        computable worlds consistent with $\rec$ that $L$ identifies after $\rec$. If
        $L$ is computable (deterministic), some world consistent with $\rec$ lies
        outside $K_L$; and if $L$ identifies after $\rec$ the eventually-zero worlds
        consistent with $\rec$, some world lies inside $K_L$.
\end{enumerate}
\end{lemma}

\begin{proof}
(i) The first claim is immediate from the definition, since $L(\rec)\in\Cons(\rec)$ has
cost $0$ and every other model cost $1$. For the second, $\Cons(\rec)$ is infinite by
\cref{thm:ud-infinite}, so it contains some $e'\ne L(\rec)$. Let $L'$ agree with $L$ on
every record other than $\rec$, and let $L'(\rec)=e'$. Then $L'$ outputs consistent
models, it differs from $L$ only on one input, so it is computable if $L$ is, and
$\Sel_{c'}(\rec)=\{e'\}\neq\{L(\rec)\}$, where $c'$ is the constraint of $L'$ given by
(i). So the record does not determine the
constraint, which is what it means for the constraint to be imposed
(\cref{def:constraint}). For (ii), the first claim is \cref{thm:ud-gold}; for the
second, the world $\sigma_\rec0^\infty$ is eventually zero and consistent with $\rec$.
\end{proof}

So choosing a method is choosing a class $K_L$ of worlds to be right about, and no
record decides whether the actual world belongs to that class. \Cref{sec:simplicity}
examines simplicity as a way of making this choice.

\paragraph{The ideal-agent concession.}
\Cref{thm:ud-gold} binds computable methods, and the restriction is not idle. A
method with an oracle for the totality problem can output the least total index
consistent with the current record; every smaller total index computes a different
function and is eventually refuted, so this \emph{identification by enumeration}
converges to a correct model on every total computable world
\citep{putnam1963,gold1967}. An ideal agent of this
kind therefore does converge on the true behavior of any deterministic computable
world, though it never knows when it has. The limit of \cref{thm:ud-gold} is
computational.
